\documentclass[9pt,aspectratio=169,professionalfonts]{beamer}
\usetheme{Madrid}
\usecolortheme{default}
\setbeamertemplate{navigation symbols}{}

% Encodage et langue
\usepackage[T1]{fontenc}
\usepackage[french]{babel}

% Packages usuels
\usepackage{amsmath,amssymb}
\usepackage{graphicx}
\usepackage{booktabs}
\usepackage{tikz}
\usepackage{listings}

% Style listings
\lstset{language=Python,basicstyle=\ttfamily\footnotesize,showstringspaces=false,breaklines=true,frame=single}

% Commandes utiles
\newcommand{\checkmarkfr}{\ding{51}}

% Titre
\title[Projection aléatoire]{Projection aléatoire : embeddings efficaces pour graphes hétérogènes}
\subtitle{Présentation technique de 30 minutes}
\author{Nom Prénom}
\institute{Institution}
\date{\today}

\begin{document}

%———————————————————————————————————
\begin{frame}
  \titlepage
\end{frame}

\begin{frame}{Plan}
  \tableofcontents
\end{frame}

%———————————————————————————————————
\section{Contexte et motivation}

\begin{frame}{Défi : réseau bibliographique}
\begin{block}{Jeu de données}
\begin{itemize}
  \item 28\,702 auteurs
  \item 20 conférences
  \item 8\,920 termes
  \item Types de relations : Auteur–Auteur, Auteur–Conférence, Auteur–Terme
\end{itemize}
\end{block}
\begin{block}{Objectif}
Apprendre des embeddings de dimension 256 préservant la sémantique afin de prédire de nouvelles collaborations.
\end{block}
\end{frame}

%———————————————————————————————————
\section{Limites des marches aléatoires}

\begin{frame}{Méthode classique : Node2Vec / MetaPath2Vec}
\begin{block}{Algorithme}
\begin{itemize}
  \item Pour chaque nœud : générer des marches aléatoires ("phrases")
  \item Entraîner Word2Vec (Skip-gram) sur ces séquences
\end{itemize}
\end{block}
\begin{alertblock}{Problèmes}
\begin{itemize}
  \item Coût computationnel et mémoire élevés
  \item Biais d'échantillonnage (structure globale manquante)
  \item Sensibilité aux hyperparamètres
\end{itemize}
\end{alertblock}
\end{frame}

%———————————————————————————————————
\section{Projection aléatoire : intuition}

\begin{frame}{Idée principale}
\begin{enumerate}
  \item Créer une matrice aléatoire creuse $\mathbf R$
  \item Appliquer directement les puissances de la matrice de voisinage $\mathbf M$ : $\mathbf M\mathbf R,\;\mathbf M^2\mathbf R,\;\mathbf M^3\mathbf R$
  \item Combiner ces vues avec des poids apprenables (softmax)
\end{enumerate}
\end{frame}

%———————————————————————————————————
\section{Fondements théoriques}

\begin{frame}{Lemme de Johnson–Lindenstrauss}
\begin{theorem}
Toute distribution aléatoire satisfaisant les hypothèses du lemme préserve les distances avec un facteur $(1\pm\varepsilon)$ après projection dans $O(\log n)$ dimensions.
\end{theorem}
\end{frame}

%———————————————————————————————————
\section{Algorithme complet}

\begin{frame}{Étapes}
\begin{enumerate}
 \item Projection aléatoire pondérée (paramètre $\alpha$)
 \item Construction / normalisation des méta-chemins (paramètre $\beta$)
 \item Calcul itératif des puissances $\mathbf U_1,\mathbf U_2,\mathbf U_3$
 \item Combinaison linéaire \(w_f\) $\rightarrow$ embedding final
 \item Perte : BCE + régularisation entropique
\end{enumerate}
\end{frame}

%———————————————————————————————————
\section{Optimisation}

\begin{frame}{Fonction de perte}
\[
\mathcal L = \mathrm{BCE}(\gamma-\lambda\|\mathbf e_i-\mathbf e_j\|^2, y_{ij}) + \lambda_{\mathrm{entropy}}\sum_f w_f\log w_f
\]
Gradients calculés par rétro-propagation ; mise à jour de $\theta_f,\lambda,\gamma$ avec Adam.
\end{frame}

%———————————————————————————————————
\section{Résultats expérimentaux}

\begin{frame}{Performances}
\begin{itemize}
  \item AUC validation : 0,94
  \item 6 × plus rapide que Node2Vec
  \item Interprétable : les poids montrent l'importance des conférences (ACA)
\end{itemize}
\end{frame}

%———————————————————————————————————
\section{Conclusion}

\begin{frame}{Points essentiels}
\begin{itemize}
  \item Projection aléatoire + puissances de matrice = embeddings efficaces
  \item Théorie garantie par Johnson–Lindenstrauss
  \item Scalabilité linéaire et support naturel des graphes hétérogènes
\end{itemize}
\end{frame}

%———————————————————————————————————
\begin{frame}{Questions ?}
\centering
\Huge Merci de votre attention !
\end{frame}

\end{document} 